% !TeX TS-program = xelatex
\PassOptionsToClass{fontset=fandol}{ctexart}
\documentclass{resume}
\ResumeName{陈沫樵}
\usepackage{enumitem}

\begin{document}

\ResumeContacts{
  (+86) 186\-1195\-1800,%
  \ResumeUrl{mailto:mochiaochen@gmail.com}{mochiaochen@gmail.com},%
  \ResumeUrl{https://github.com/MochiaoChen}{github.com/MochiaoChen}
}

\ResumeTitle

\section{学校}
\ResumeItem
[对外经济贸易大学 | 本科]
{\textbf{对外经济贸易大学}}
[{\zihao{5} 金融学（鸿儒金融人才培养实验班），中国金融学院 | 经济学学士}]
[2024.09—2028.06]

\section{实习}

\ResumeItem{\textbf{AI 应用与评测实习生}}
[UniPat AI]
[2026.08—2026.09]
\begin{itemize}
  \item \textbf{AI 质检产品设计}：针对专家审核中条件与示例错配、原文查找困难、反馈难以回溯等问题，梳理“评价问题—逻辑关系—检查条件”的数据结构，推进\textbf{审核状态、修改建议、结果导出及 PDF 证据回溯}原型；结合实际案例开展字段校验与交互验收，形成可复用的专家复核流程。
  \item \textbf{垂域评测体系建设}：面向美股公司研究中的信息理解、财务推导与复杂推理任务，参与构建结构化评分标准（\textbf{Rubric}），将专家分析拆解为可核验的事实、逻辑关系与通过条件；推进代表性案例由 \textbf{59} 项向 \textbf{77} 项原子化标准迭代，明确证据要求、评分权重与重复扣分边界。
  \item \textbf{模型效果分析与优化}：结合 \textbf{LLM-as-a-Judge}、程序化财务检查和多轮重复实验开展 \textbf{Failure Case} 分析，区分内容遗漏、证据提取失败、口径冲突与评审误判；定位指标归一化、跨期匹配和高权重扣分误触发等问题，为数据规则、Prompt 与评测流程迭代提供依据。
  \item \textbf{AI 工作流与质检工具}：使用 AI 编码工具辅助 \textbf{Python}、\textbf{JavaScript} 开发，推进结构化材料生成与专家审核原型，支持检查条件与示例对齐、审核状态管理、修改反馈及结果导出；完成单案例 \textbf{300} 余项审核材料的生成与结构校验，相关工具版本通过 \textbf{36} 项自动化测试。
  \item \textbf{数据治理与证据追溯}：参与建设 \textbf{1,700} 余份研报的资料储备，完成元数据整理、分类映射、附件核对与异常记录；推进“评价条目—基础证据—原始 PDF”的追溯流程，支持正文、表格与跨页图表候选定位、页面高亮和人工复核，覆盖单案例 \textbf{114} 条证据锚点。
\end{itemize}

\ResumeItem{\textbf{产品经理实习生}}
[Tara]
[2026.07]
\begin{itemize}
  \item \textbf{产品设计}：负责 \textbf{AI 社交产品}用户端功能设计，梳理用户需求与核心使用流程，撰写 \textbf{PRD} 并输出原型与交互说明；完成用户卡片体系（区分供给侧从业者与需求侧学生用户）、核心入口命名与信息架构设计；基于社会学 \textbf{Recognition} 理论与“同好社区”vs“社交网络”的定位差异，提炼具体产品设计原则。
  \item \textbf{原型开发}：独立使用 \textbf{AI 辅助编程}工具完成\textbf{微信小程序}前端交互开发，承担从原型设计到功能实现的全流程，推动产品从概念快速落地为可交互原型。
  \item \textbf{市场活动策划}：策划并执行面向北京科技公司实习生群体的 \textbf{Skill Chat} 线下活动，负责活动定位、目标人群筛选与内容形式设计，探索通过垂直兴趣社群实现精准获客的运营路径。
  \item \textbf{市场分析与 BD}：持续跟踪 AI 社交赛道竞品动态与产品机制（包括 A2A 社交、慧眸等相关产品案例），为公司定位与差异化策略提供分析支持；对接潜在合作方与资源方，推进早期商务合作洽谈。
\end{itemize}

\ResumeItem{\textbf{投资者关系实习生}}
[奇绩创坛]
[2026.04—2026.06]
\begin{itemize}
  \item \textbf{投研支持}：研究 \textbf{Agentic AI} 领域发展态势，为陆奇博士的研究提供支持；对市场进行 mapping，梳理潜在投资机会。
  \item \textbf{创业者网络}：链接 \textbf{500+} 创业者，对其学术禀赋与产业背景进行综合评估，定期跟踪动态并协助市场研究与数据分析，系统性研究潜在投资标的，及时发现新兴投资机会。
  \item \textbf{生态拓展}：链接 \textbf{50+} 海内外高校研究人员，针对视频生成、Agentic AI 等前沿领域通过深度访谈、行业调研、桌面研究快速开展研究，挖掘算力需求方并提供算力支持；直接对接高校科研团队、技术创业者及产业专家，通过社区网络、高校演讲等渠道挖掘潜力项目与人才。
  \item \textbf{项目管理}：制定活动方案并组织执行，调研分析创业者真实需求与反馈，协助平台资源对接及被投企业需求响应。
\end{itemize}

\ResumeItem{\textbf{投资银行部实习生}}
[中信证券]
[2026.01—2026.04]
\begin{itemize}
  \item \textbf{行业洞察}：深度参与某具身智能企业港股 \textbf{IPO} 项目，参与招股书中“业务”章节的撰写。通过对 AI 算法、感知系统及机器人硬件链条的梳理，完成核心业务逻辑分析与多家同业竞对研究。
  \item \textbf{流程支持}：负责协助 \textbf{M111} 表格填报及港交所回函撰写，确保上市申请流程符合港交所合规标准。
  \item \textbf{AI 提效}：使用 \textbf{Claude Code} 完成工作中大量重复性任务（如底稿核对、格式排版、资料整理等），显著提升项目组交付效率。
\end{itemize}

\ResumeItem{\textbf{投资银行部实习生}}
[长城证券]
[2025.09—2026.03]
\begin{itemize}
  \item \textbf{项目承做}：深度参与\textbf{公司债}项目的\textbf{尽职调查}，协助撰写立项报告与申报材料，整理底稿并完成归档。
  \item \textbf{数据分析}：运用 \textbf{Wind}、\textbf{iFind}、\textbf{Excel} 及 \textbf{Python} 进行市场数据清洗与同行业分析。
  \item \textbf{Vibe Coding 提效}：使用 \textbf{Vibe Coding} 快速开发小工具，辅助完成数据整理、材料核对等重复性、机械性工作，降低人工操作成本。
\end{itemize}

\ResumeItem{\textbf{AI 科技内容运营实习生}}
[云杉资本]
[2025.03—2025.09]
\begin{itemize}
  \item \textbf{行业研究}：聚焦 \textbf{AI} 与\textbf{消费}赛道，独立撰写多篇行业分析简报，拆解国内外 AI 初创项目的商业模式与技术路径。
  \item \textbf{内容策划}：建立 AI 赛道投资视角话题库，策划并输出深度分析文章，辅助投资团队进行市场趋势研判。
\end{itemize}

\section{项目}

\ResumeItem{\textbf{鉴真盾}}
[Qwen3.6、Agent、Prompt Engineering、Fine-tuning]
[2025.11]
\begin{itemize}
  \item \textbf{项目概况}：开发“鉴真盾”——基于\textbf{多模态大模型}的金融反欺诈人像风险识别系统，面向小微企业贷款预审场景，解决黑产中介伪造身份、批量攻击等传统信审痛点，通过人像识别判定申请人真实性风险并输出量化评分。
  \item \textbf{技术实现}：基于 \textbf{Qwen3.6} 模型完成数据标注、\textbf{Agent} 协作架构设计、模型测试与 \textbf{Prompt Engineering} 调优，并对模型进行 \textbf{Fine-tune}；建立中介欺诈特征库，精准识别六大类欺诈特征，负责全栈开发。
  \item \textbf{量化成果}：模型在测试集准确率达 \textbf{92\%}，实现风险量化评分机制；项目获“三创赛”\textbf{北京市一等奖}、“工行杯”全国大学生金融科技创新大赛\textbf{全国三等奖}。
\end{itemize}

\ResumeItem{\textbf{Lectio}}
[Agent Skill、\LaTeX{}、GLM]
[2025.11]
\begin{itemize}
  \item \textbf{项目概况}：独立设计并开发开源 AI 学习 Skill \textbf{Lectio}，可将 PDF 课件、B 站/YouTube 视频及零散笔记一键重构为结构化 \textbf{\LaTeX{}} 讲义、Markdown 预览及 Anki 闪卡，面向 agent 运行时设计，在智谱 \textbf{GLM} 系列模型上调优。
  \item \textbf{产品设计}：内置苏格拉底思考题、费曼压缩、自动闪卡等教学法增强模块；项目以 \textbf{GPL-3.0} 协议开源，获北京市大学生计算机设计大赛\textbf{北京市三等奖}。
\end{itemize}

\ResumeItem{\textbf{Job Seeking——端到端 AI 求职助手}}
[Agent Skill、Python、Gmail]
[2026.07]
\GrayText{项目地址：\ResumeUrl{https://github.com/MochiaoChen/interview-prep-skill}{github.com/MochiaoChen/interview-prep-skill}}
\begin{itemize}
  \item \textbf{产品定义}：针对求职过程中简历、JD、求职信、邮件投递与面试准备相互割裂的问题，从 0 到 1 设计覆盖完整求职周期的 \textbf{AI Agent}，让用户通过自然语言完成材料管理、岗位匹配、投递和面试准备。
  \item \textbf{核心机制}：以个人简历作为统一信息源，根据不同 JD 自动提取岗位要求、匹配候选人经历并生成针对性求职信，减少重复填写和多工具切换，同时避免 AI 编造个人经历。
  \item \textbf{数据沉淀}：将每次申请的公司、岗位、JD、求职信和简历附件沉淀为独立档案，并自动复用于后续面试问题预测和准备手册生成，解决投递信息难以持续利用的问题。
  \item \textbf{交互设计}：设计“默认生成邮件草稿、确认后再发送”的关键交互，在提升投递效率的同时降低错发、附件错误和隐私泄露风险；使用 \textbf{Agent Skill}、\textbf{Python} 与 \textbf{Gmail} 完成产品落地。
\end{itemize}

\ResumeItem{\textbf{Text2Inflation——基于央行文本叙事的通胀预测研究}}
[{\zihao{5} Python、LLM、NLP、机器学习、计量经济学}]
[2026]
\begin{itemize}
  \item \textbf{项目概况}：独立完成研究与开发，运用 \textbf{Python}、\textbf{LLM}、\textbf{NLP}、机器学习及计量经济学方法，构建端到端通胀预测框架，利用 LLM 解析中国人民银行《货币政策执行报告》，将非结构化政策文本转化为通胀情绪、政策立场、风险归因等 \textbf{10 维}叙事特征；论文获鸿儒本科生学术论坛\textbf{论文三等奖}。
  \item \textbf{模型体系}：实现 LASSO、Elastic Net、随机森林、PCA、PLS、LSTM、XGBoost 等基准与文本增强模型，通过单步滚动扩展窗口预测，避免随机划分造成的时间序列信息泄漏。
  \item \textbf{评估体系}：建立完整的模型评估体系，使用 RMSE、MAE、$R^2$ 比较预测效果，并结合 Diebold–Mariano 检验、SHAP 特征归因及交互效应分析，验证文本信息的增量预测价值。
  \item \textbf{稳健性分析}：针对小样本和宏观经济结构变化，进一步开展状态条件检验、发布时滞敏感性分析及通缩/大宗商品冲击案例研究，分析央行叙事在不同经济阶段中的作用机制。
  \item \textbf{工程实现}：将 PDF 解析、LLM 特征提取、模型训练、预测评估和可解释性分析整合为一键式流水线，实现研究过程的自动化与可复现；项目地址：\ResumeUrl{https://github.com/MochiaoChen/Text2Inflation}{github.com/MochiaoChen/Text2Inflation}。
\end{itemize}

\ResumeItem{\textbf{NIPT 时点优化与异常判定研究}}
[Python、机器学习、数据可视化]
[2025.10]
\begin{itemize}
  \item \textbf{项目概况}：针对无创产前检测（NIPT）时点优化与异常判定问题，基于临床数据构建多维度机器学习模型体系。
  \item \textbf{数据分析与建模}：负责完成原始数据缺失值插补、异常值剔除等预处理工作；协助建模手完成编程工作。
  \item \textbf{成果产出}：负责撰写项目核心摘要，并使用 Python 绘制数据分布、特征重要性、ROC 曲线等多维可视化图表；项目获全国大学生数学建模竞赛\textbf{北京赛区一等奖}。
\end{itemize}

\section{校园组织经历}

\ResumeItem{\textbf{UIBE AI 创协}}
[Co-founder]
[2025.11—至今]
\begin{itemize}
  \item \textbf{组织搭建}：从 0 到 1 发起 AI 创协，搭建校园运营、Startup 合作与 AI 应用三大板块，沉淀 \textbf{200+} 位活跃社群用户。
  \item \textbf{SOP 构建}：面向 AI 企业设计校园增长合作模式 SOP，通过“校园运营—定制服务—数据复盘”的闭环，为企业提供校园增长解决方案；合作企业与生态覆盖智谱、Second Me、PandaAI、HyperKnow、观猹、知乎等。
  \item \textbf{营销实验}：调研学生群体用户需求并制定运营方案，通过线下活动、社媒、社群等多渠道投放与互动，使用 AI 优化流程（如小红书自动图文 Agent）并根据数据调整运营策略，两周内累计为合作企业获取用户 \textbf{200+}。
\end{itemize}

\section{能力}
\begin{itemize}
  \item \textbf{AI 工具与开发范式}：熟练使用 \textbf{Codex} 与 \textbf{Claude Code}（三个月 token 耗量近 \textbf{60 亿}）进行 \textbf{Vibe Coding} 开发；具备较强学习能力与 \textbf{AI Native} 的思考与协作方式，能快速上手并驾驭新兴 AI 开发工具链。
  \item \textbf{编程与数据分析}：熟练掌握 \textbf{Python}，擅长数据分析；了解 \textbf{SQL}、\textbf{C++}、\textbf{R}，对现代软件开发流程有基本认知；掌握机器学习基础与数理统计建模，能够独立完成数据清洗、可视化分析及行业研究简报撰写。
  \item \textbf{内容与网感}：具备较强网感，独立运营小红书账号——单周浏览量 \textbf{37.6w}，封面点击率 \textbf{23.7\%}，获赞 \textbf{3932}、评论 \textbf{3446}、收藏 \textbf{1785}、分享 \textbf{3086}，涨粉 \textbf{643}；产出 \textbf{10w+} 爆款 1 篇、\textbf{5w+} 2 篇、\textbf{1w+} 10 篇；对新领域、新内容具备独立思考能力与批判性思考能力。
  \item \textbf{语言能力}：英语 \textbf{C1} 水平（CET-4 \textbf{626} 分），具备优秀的英文专业文献阅读与学术写作能力，英语可作为工作语言。
\end{itemize}

\end{document}
